\section{Исходная программа}

Сборка и запуск:
\begin{itemize}
  \item Windows: \texttt{cmake -S . -B build-windows \&\& cmake --build build-windows}, запуск
    \texttt{./build-windows/bin/Release/parent.exe} из папки с \texttt{child.exe};
  \item Linux: \texttt{cmake -S . -B build-linux \&\& cmake --build build-linux}, запуск
    \texttt{./build-linux/bin/parent} из папки с \texttt{child};
  \item файл с исходными текстами содержит папки \texttt{src}, \texttt{app} и корневой
    \texttt{CMakeLists.txt}.
\end{itemize}

Единый интерфейс:

\begin{center}
\lstinputlisting[language=C,caption=src/os\_api.h,captionpos=b]{inc/os_api.h}
\end{center}

Реализация для Windows:

\begin{center}
\lstinputlisting[language=C,caption=src/os\_api\_win.c,captionpos=b]{inc/os_api_win.c}
\end{center}

Реализация для Linux:

\begin{center}
\lstinputlisting[language=C,caption=src/os\_api\_posix.c,captionpos=b]{inc/os_api_posix.c}
\end{center}

Родительский процесс:

\begin{center}
\lstinputlisting[language=C,caption=app/parent.c,captionpos=b]{inc/parent.c}
\end{center}

Дочерний процесс:

\begin{center}
\lstinputlisting[language=C,caption=app/child.c,captionpos=b]{inc/child.c}
\end{center}

Файл сборки:

\begin{center}
\lstinputlisting[language=C,caption=CMakeLists.txt,captionpos=b]{inc/CMakeLists.txt}
\end{center}

Ниже приведены ключевые системные вызовы родительского процесса под Linux (полная трасса ---
\texttt{strace -f -o strace.txt ./parent}, итог сохранён в \texttt{strace.txt}):
\begin{lstlisting}[language=C,caption=strace (fragment),captionpos=b]
execve("./parent", ["./parent"], ...) = 0
pipe2([3, 4], 0) = 0                          // pipe: in  = [3, 4]
pipe2([5, 6], 0) = 0                          // pipe: out = [5, 6]
rt_sigaction(SIGPIPE, {sa_handler=SIG_IGN, ...}, ...) = 0
clone(child_stack=NULL, flags=...|SIGCHLD, ...) = 1303 // fork()
close(3) = 0                                  // parent: in.read
close(6) = 0                                  // parent: out.write
  dup2(3, 0) = 0                              // child: stdin  <- in.read
  dup2(6, 1) = 1                              // child: stdout <- out.write
  close(4) = 0
  close(5) = 0
  execve("./child", ["child", "test_strace"], ...) = 0
write(4, <1.0f = 0x3F800000>, 4) = 4          // parent: float 1.0 -> in.write
read(3, <1.0f = 0x3F800000>, 4096) = 4        // child:  reads 1.0
write(4, <2.0f = 0x40000000>, 4) = 4          // parent: float 2.0 -> in.write
read(3, <2.0f = 0x40000000>, 4096) = 4        // child:  reads 2.0
write(4, <3.0f = 0x40400000>, 4) = 4          // parent: float 3.0 -> in.write
read(3, <3.0f = 0x40400000>, 4096) = 4        // child:  reads 3.0
write(4, <0.0f = 0x00000000>, 4) = 4          // parent: float 0.0 -> in.write
read(3, <0.0f = 0x00000000>, 4096) = 4        // child:  reads 0.0
close(4) = 0                                  // parent: EOF for child
read(3, "", 4096) = 0                         // child:  sees EOF
write(1, <6.0f = 0x40C00000>, 4) = 4          // child:  sum 6.0 -> stdout
read(5, <6.0f = 0x40C00000>, 4) = 4           // parent: reads sum
wait4(1303, [{WIFEXITED(s) && WEXITSTATUS(s) == 0}], 0, NULL) = 1303
exit_group(0) = ?
\end{lstlisting}
Трассировка подтверждает описанный ход работы: создание двух каналов (\texttt{pipe2}), подавление
\texttt{SIGPIPE} (\texttt{rt\_sigaction}), запуск дочернего процесса (\texttt{clone}), подключение
каналов к стандартному вводу/выводу (\texttt{dup2}) с последующим \texttt{execve("./child")},
передачу чисел и итоговой суммы (4 байта в младшей части типа \texttt{float}), корректную передачу
\texttt{EOF} путём закрытия записывающего конца (\texttt{close(4)} $\rightarrow$ \texttt{read = 0})
и ожидание завершения процесса (\texttt{wait4}) с кодом выхода \texttt{0}.